\documentclass[11pt,spanish]{article}

% =========================================================
% PLANTILLA BASE PARA INFORMES ACADÉMICOS / TÉCNICOS
% =========================================================

% --- Idioma y codificación ---
\usepackage[spanish,es-tabla]{babel}
\usepackage[T1]{fontenc}
\usepackage{textcomp}
\usepackage{newunicodechar}
\newunicodechar{₡}{\textcolonmonetary}

% --- Márgenes / papel ---
\usepackage{geometry}
\geometry{a4paper, margin=2.7cm}

% --- Matemáticas ---
\usepackage{amsmath}

% --- Tablas, enlaces y elementos visuales ---
\usepackage{array}
\usepackage{tabularx}
\usepackage{booktabs}
\usepackage{longtable}
\usepackage{graphicx}
\usepackage{float}
\usepackage{xcolor}
\usepackage{tikz}
\usetikzlibrary{positioning,shapes.geometric,arrows.meta,calc}
\usepackage{enumitem}
\usepackage{ragged2e}
\usepackage[hidelinks]{hyperref}

% --- Encabezado y pie de página ---
\usepackage[myheadings]{fancyhdr}
\usepackage{lastpage}
\setlength{\headheight}{30pt}
\setlength{\footskip}{30pt}

% --- Interlineado ---
\linespread{1.2}
\sloppy

% =========================================================
% DATOS DEL DOCUMENTO
% =========================================================
\newcommand{\Institucion}{Instituto Tecnológico de Costa Rica}
\newcommand{\Campus}{Campus Tecnológico Central Cartago}
\newcommand{\DocumentoTitulo}{Reflexiones y Recomendaciones}
\newcommand{\DocumentoSubtitulo}{}
\newcommand{\Curso}{Administración de proyectos}
\newcommand{\Profesor}{Alicia Marcela Salazar Hernández}
\newcommand{\FechaDocumento}{23 de junio de 2026}
\newcommand{\PeriodoAcademico}{I Semestre}
\newcommand{\AnioDocumento}{2026}

\newcommand{\AutoresPortada}{%
    Mauricio González Prendas 2024143009\\
    Isaac Jiménez Blanco 2024202804\\
    Tamara Robles Camacho 2024099342%
}

% Datos que aparecen en el encabezado.
\newcommand{\DocHeaderLeft}{}
\newcommand{\DocHeaderRight}{Reflexiones y Recomendaciones}
\newcommand{\DocFooterRight}{Página~\thepage\ de~\pageref{LastPage}}

% Metadatos PDF.
\title{\DocumentoTitulo}
\author{\AutoresPortada}
\date{\FechaDocumento}

% =========================================================
% CONFIGURACIÓN DE TABLAS Y UTILIDADES
% =========================================================
\newcolumntype{Y}{>{\RaggedRight\arraybackslash}X}
\renewcommand{\arraystretch}{1.22}
\setlength{\LTpre}{0.5em}
\setlength{\LTpost}{0.5em}

% Imagen segura si el archivo no existe, muestra un recuadro de marcador.
\newcommand{\SafeImage}[2][0.92\textwidth]{%
    \IfFileExists{#2}{%
        \includegraphics[width=#1]{#2}%
    }{%
        \fbox{%
            \begin{minipage}[c][4cm][c]{#1}
            \centering
            Imagen pendiente\\
            \texttt{#2}
            \end{minipage}%
        }%
    }%
}

% Figura reutilizable.
\newcommand{\EvidenceFigure}[3]{%
    \begin{figure}[H]
    \centering
    \SafeImage{#1}
    \caption{#2}
    \label{#3}
    \end{figure}%
}

% =========================================================
% ENCABEZADO Y PIE
% =========================================================
\pagestyle{fancy}
\fancyhf{}
\fancyhead[L]{\DocHeaderLeft}
\fancyhead[R]{\DocHeaderRight}
\renewcommand{\headrulewidth}{0.4pt}
\renewcommand{\footrulewidth}{0.4pt}
\fancyfoot[R]{\DocFooterRight}

% Estilo equivalente para páginas horizontales.
\fancypagestyle{widefancy}{%
    \fancyhf{}%
    \fancyhead[L]{\DocHeaderLeft}%
    \fancyhead[R]{\DocHeaderRight}%
    \renewcommand{\headrulewidth}{0.4pt}%
    \renewcommand{\footrulewidth}{0.4pt}%
    \fancyfoot[R]{\DocFooterRight}%
}

% =========================================================
% PÁGINAS HORIZONTALES PARA TABLAS ANCHAS
% =========================================================
\makeatletter
\newcommand{\SyncPageBuilderDimensions}{%
    \global\vsize=\textheight
    \global\@colroom=\textheight
    \global\@colht=\textheight
}
\makeatother

\newenvironment{widepage}{%
    \clearpage
    \hypersetup{pageanchor=false}%
    \paperwidth=297mm
    \paperheight=210mm
    \pdfpagewidth=297mm
    \pdfpageheight=210mm
    \newgeometry{left=2.7cm,right=2.7cm,top=2.7cm,bottom=2.7cm}%
    \setlength{\textwidth}{243mm}%
    \setlength{\textheight}{156mm}%
    \setlength{\linewidth}{\textwidth}%
    \setlength{\columnwidth}{\textwidth}%
    \hsize=\textwidth
    \setlength{\headwidth}{\textwidth}%
    \SyncPageBuilderDimensions
    \pagestyle{widefancy}%
}{%
    \clearpage
    \paperwidth=210mm
    \paperheight=297mm
    \pdfpagewidth=210mm
    \pdfpageheight=297mm
    \restoregeometry
    \SyncPageBuilderDimensions
    \hypersetup{pageanchor=true}%
    \pagestyle{fancy}%
}

% =========================================================
% PORTADA
% =========================================================
\newcommand{\MakeCover}{%
\begin{titlepage}
\thispagestyle{empty}
\centering

{\bfseries \huge \Institucion \par}
\vspace{0.25cm}
{\LARGE \Campus \par}

\vfill

{\huge \DocumentoTitulo \par}

\vfill

{\bfseries \LARGE Profesora \par}
{\Large \Profesor \par}

\vfill

{\bfseries \LARGE Estudiantes \par}
{\Large \AutoresPortada \par}

\vfill

{\bfseries\LARGE Fecha \par}
{\Large \FechaDocumento \par}

\vfill

{\LARGE \PeriodoAcademico \par}
{\LARGE \AnioDocumento \par}

\end{titlepage}%
}

% =========================================================
% DOCUMENTO
% =========================================================
\begin{document}
\hypersetup{pageanchor=false}

\MakeCover

\pagenumbering{roman}
\pagestyle{empty}
\tableofcontents
\clearpage
\pagenumbering{arabic}
\hypersetup{pageanchor=true}
\pagestyle{fancy}

\section{Información general del documento}

\begin{table}[H]
\centering
\begin{tabularx}{\textwidth}{p{5cm} Y}
\toprule
\textbf{Nombre del proyecto} & Plataforma Universitaria Integrada \\
\textbf{Organización} & Universidad Tecnológica La Mejor \\
\textbf{Documento} & \DocumentoTitulo \\
\textbf{Project manager} & Tamara Robles Camacho \\
\textbf{Fecha} & \FechaDocumento \\
\bottomrule
\end{tabularx}
\end{table}

\section{Reflexión sobre la gestión del proyecto}

La gestión del proyecto se caracterizó por una aplicación rigurosa de los grupos de procesos propuestos por el estándar en medio de un marco de tiempo sumamente desafiante debido a que el trabajo inició el 8 de junio y debía entregarse el 24 de junio de 2026. Durante la fase de inicio y planificación se logró establecer un alcance realista y una estructura de trabajo bien delimitada ya que los documentos como el acta de constitución y el cronograma brindaron el horizonte necesario para ordenar las tareas de todo el equipo de desarrollo. Asimismo la ejecución y el monitoreo permitieron identificar de manera oportuna las desviaciones en el esfuerzo de desarrollo del Front-End por lo tanto la implementación del control de cambios funcionó como una herramienta estratégica para proteger la calidad del producto y asegurar que el proyecto finalizara exitosamente en la fecha estipulada por la institución.

\section{Reflexión sobre el trabajo en equipo}

El desempeño colaborativo del grupo compuesto por cinco miembros resultó ser uno de los pilares más fuertes de la iniciativa ya que la división clara de los roles entre documentadores y desarrolladores permitió avanzar en paralelo sin generar cuellos de botella importantes. De esta manera el uso de herramientas de comunicación inmediata como WhatsApp agilizó las consultas cotidianas y los acuerdos rápidos mientras que el uso de GitHub garantizó un control efectivo de las versiones del código del prototipo web en cada fase. En consecuencia las reuniones de seguimiento regulares facilitaron la sincronización entre el diseño técnico y los artefactos de gestión para lo cual el liderazgo de la administración del proyecto logró integrar las diferentes perspectivas y mantener la motivación a pesar de la presión fuerte del cronograma semanal.

\section{Reflexión sobre la gestión de riesgos}

El análisis de los 16 riesgos identificados en la planificación demostró que el equipo anticipó correctamente los principales retos que afectaban el proyecto ya que la restricción de tiempo y la complejidad técnica figuraron como las amenazas de mayor impacto desde el primer momento de trabajo. Por lo tanto algunos riesgos críticos como el incremento desmedido del alcance y la subestimación del esfuerzo de desarrollo se materializaron parcialmente cuando se solicitaron funcionalidades adicionales en el portal estudiantil durante el ciclo iterativo. Esto implica que las estrategias de mitigación establecidas fueron completamente efectivas debido a que la reducción del alcance en el dashboard ejecutivo permitió controlar el riesgo R-04 y el riesgo R-10 garantizando así que el producto conservara su estabilidad general y se entregara sin defectos mayores a los usuarios finales.

\section{Reflexión sobre la calidad del proyecto}

El cumplimiento de los estándares de calidad se evidenció a través de las diversas revisiones aplicadas tanto a la documentación como al prototipo final debido a que el grupo se apegó a los criterios de aceptación definidos para cada módulo de la plataforma universitaria. Asimismo las pruebas manuales ejecutadas sobre el Front-End aseguraron que la navegación y el diseño visual ofrecieran una experiencia intuitiva para los estudiantes docentes y administrativos ya que se priorizaron los elementos funcionales sobre los adornos innecesarios en la interfaz gráfica. Además de esto la revisión cruzada de los archivos escritos garantizó una alta coherencia entre los presupuestos los cronogramas y los planes de gestión por lo tanto el resultado refleja un esfuerzo sistemático por entregar un trabajo académico profesional confiable y sumamente ordenado.

\section{Recomendaciones para proyectos futuros}

A partir de las lecciones asimiladas durante la ejecución se recomienda definir el alcance del producto con el mayor nivel de detalle posible en las primeras reuniones de inicio ya que esto ayuda a prevenir discusiones complejas cuando el tiempo de trabajo se encuentra avanzado. Asimismo resulta fundamental asignar un margen de contingencia amplio en el cronograma para todas las tareas de programación y revisión técnica debido a que el desarrollo de interfaces suele enfrentar imprevistos que consumen más horas de las planificadas originalmente en las estimaciones. Además de estas medidas se sugiere mantener una plantilla unificada para todos los documentos del proyecto desde el día uno por lo tanto el equipo evita retrabajos de formato continuos y asegura que la imagen del trabajo sea consistente ante los evaluadores académicos.

Por otro lado es recomendable establecer canales formales de aprobación de cambios que funcionen de manera ágil y sin excesiva burocracia ya que en proyectos cortos las decisiones deben tomarse rápidamente para no paralizar el avance de los programadores front-end. Finalmente se aconseja fomentar una cultura de revisión por pares en todos los entregables técnicos y documentales para lo cual cada miembro debe comprometerse a leer el trabajo de sus compañeros antes de consolidar la versión definitiva mitigando así los errores por omisión y elevando sustancialmente el estándar general del proyecto universitario.

\section{Conclusión general del proyecto}

El desarrollo de la plataforma universitaria integrada representa un logro valioso que combina la disciplina de la gestión de proyectos con las habilidades técnicas de la construcción de sistemas web debido a que el equipo consiguió transformar una necesidad institucional abstracta en un prototipo funcional. De esta manera el esfuerzo invertido en cada fase demuestra que la planificación adecuada y el control de los recursos son herramientas indispensables para superar las barreras del tiempo estricto y las limitaciones tecnológicas en cualquier entorno de trabajo colaborativo. En consecuencia la experiencia adquirida durante este corto período fortalece las competencias profesionales de todos los participantes dedicados y establece una base sólida para enfrentar futuros desafíos en el desarrollo de software y la administración experta de requerimientos complejos.

\end{document}
